\section{AGI 路线图：从 ChatGPT 山脚到行动主峰}

前面几章分别讲模型上限、行动环境和公司分化，本章进入广密作为“AGI 原教旨主义者”的路线图。他认为智能提升是唯一主线，智能本身就是最大应用。ChatGPT 只是山脚，后面还有 Coding、Coding Agent、General Agent、AI for Science、Robotics 等多个山头。

\begin{figure}[H]
\centering
\includegraphics[width=0.96\textwidth]{agi-roadmap-mountain.png}
\caption{AGI 路线山脊：ChatGPT 只是山脚，后续主峰依次进入更强行动环境。自制概念图，依据 00:30:18--00:43:00 对谈内容整理。}
\end{figure}

\begin{knowledgebox}{读图：主峰从对话转向行动}
从 ChatGPT 到 Coding、Coding Agent、General Agent、AI4Science、Robotics，路线越来越接近真实行动与真实反馈。聊天证明模型能表达智能；Agent 和科学证明模型能使用智能改变环境。
\end{knowledgebox}

\subsection{智能本身就是最大应用}

本节解释“智能本身就是最大应用”。这句话不是否认具体产品，而是强调所有产品都在承接模型能力外溢。ChatGPT、Cursor、Manus、Perplexity 等产品之所以出现，是因为底座模型能力突然足够强，产品公司用一个环境或容器把它变成用户可感知的工作流。应用公司不是凭空创造智能，而是在合适时间承接研究溢出的红利。

\begin{importantbox}{应用公司的任务：构建智能容器}
如果底座模型是水位上升，应用公司的任务不是假装自己造了海，而是建好港口、管道和水龙头。容器越贴近真实任务、反馈越清晰、成本越可控，越有机会沉淀壁垒。
\end{importantbox}

\subsection{AI for Science 和 Robotics 的位置}

本节看路线图后段。广密认为 2026/2027 年可能出现 AI for Science 的爆发，因为科学问题有高价值、可验证反馈和巨大市场空间。Robotics 更接近物理世界，但他对当前 robotics foundation model 的做法保持怀疑：如果只是把 VLM/VLA 套到机器人上，而没有真正解决行动、数据、反馈和本体问题，就不够本质。

\begin{warningbox}{Robotics 不应被过早神话}
机器人是 AGI 路线的重要终局，但它的数据、硬件、物理反馈和安全约束比数字世界复杂得多。相比先在 Coding 和 Agent 环境里获得稳定反馈，直接押物理世界可能会更慢、更贵、更难验证。
\end{warningbox}

\subsection{烟雾弹与主线：文生图该放在哪里}

本节处理一个容易争议的判断。广密提到文生图可能是 OpenAI 的烟雾弹，意思不是图像生成没有价值，而是它可能吸引过多注意力，却不一定代表智能主线的最高优先级。对于 AGI 原教旨主义者，图像生成、短视频、消费级爆点都要回到同一个问题：它是否推动模型更聪明、更能行动、更能学习？

\begin{knowledgebox}{课堂提示：用“是否增加行动能力”筛选 AI 热点}
一个 AI 热点可以很商业、很酷、很传播，但如果它不增加模型理解世界、操作工具、获得反馈或迁移能力的深度，它就未必是 AGI 主线。季报的价值正在于用主线过滤热点。
\end{knowledgebox}

\subsection{本章小结}

广密的 AGI roadmap 把 ChatGPT 放在山脚，把行动环境和科学发现放在后续主峰。这个路线图未必准确预测时间，但它提供了一个判断框架：越能让模型行动、验证、学习和改变世界，越靠近主线。

